\documentclass[11pt]{article}
\usepackage[margin=2.5cm]{geometry}
\usepackage{amsmath,amssymb,amsthm,mathtools}
\usepackage{booktabs,enumitem}
\usepackage[hidelinks]{hyperref}

\newtheorem{theorem}{Theorem}
\newtheorem{proposition}[theorem]{Proposition}
\newtheorem{lemma}[theorem]{Lemma}
\newtheorem{corollary}[theorem]{Corollary}
\theoremstyle{definition}
\newtheorem{definition}[theorem]{Definition}
\newtheorem{example}[theorem]{Example}
\theoremstyle{remark}
\newtheorem{remark}[theorem]{Remark}

\newcommand{\Ks}{\mathcal{K}^{\star}}
\newcommand{\thl}{\underline{\theta}}
\newcommand{\thu}{\overline{\theta}}
\newcommand{\0}{\mathbf{0}}

\title{HICA-S, Track T4 (addendum):\\ An FD-Completion Label Rule for Algorithm~A}
\author{Draft for supervisor review --- Tran An Khoa}
\date{September 2026}

\begin{document}
\maketitle

\begin{abstract}
The companion note (\emph{The Local Pruning Frontier}, hereafter [LPF]) characterises the
smallest pool $\Ks_i$ that a bid-independent, driver-local rule may retain, but computes
it only \emph{after} Algorithm~A has enumerated every route. This addendum moves part of that
reasoning \emph{inside} the label-extension procedure. When a partial route has served, or
picked up, a set $T$ of orders, we bound from below --- uniformly over all feasible
completions and all admissible bids --- how much the route would save by dropping $T$ and
sending it to the fixed fleet; if the saving is at least the FD price $q(T)$, the label and its
whole subtree are discarded. We prove that the rule never discards a route of $\Ks_i$, so that
Algorithm~A with the rule followed by $\Pi^\star$ produces exactly the same pool, allocation,
and VCG payments as before. The rule is an adaptation of \emph{rollback pruning}
(Lozano, Duque \& Medaglia, 2016); the adaptation is not a formality: because working time
enters the objective and waiting is allowed, the direct transfer of rollback is unsafe
(Example~\ref{ex:counter}), and a correction for waiting absorption
(Lemma~\ref{lem:absorb}) is required. We also report a negative result: a forward
lower-bound matrix, the most natural first idea, is too weak to ever fire on the test
instances, for a structural reason. An independent brute-force audit found no violation of any
bound, and mutation tests confirm that each component of the bound is necessary.
\end{abstract}

%=====================================================================
\section{Setting and notation}
%=====================================================================

Notation follows [LPF]. Fix a driver $i$ with start point $v_0$, departure time $t_0$ and
coefficient $\kappa$. A \emph{label} $L=(v_1,\dots,v_m)$ is a precedence-feasible stop
sequence (pickups $p_o$, deliveries $d_o$). Travel distances $d(\cdot,\cdot)$ and times
$\tau(\cdot,\cdot)$ satisfy the triangle inequality, service times $s_v\ge 0$, and a pickup
$p_o$ has an opening (release) time $e_{p_o}$ that is met by waiting; deliveries have no
opening time (Remark~\ref{rem:dlb} treats the general case). Service starts and ready times
are
\[
a_k=\max\{e_{v_k},\,r_{k-1}+\tau(v_{k-1},v_k)\},\qquad r_k=a_k+s_{v_k},\qquad r_0=t_0 .
\]
For a gigworker, $K=\kappa\cdot(\text{route length})$ and $W=(r_{\text{last}}-t_0)/60$; for an
occasional driver the route ends with a leg to the personal destination (no opening time),
$K=\kappa\cdot(\text{length}-\text{direct length})$ and
$W=(\text{arrival at destination}-t_0-\text{direct time})/60$. The reported cost is
$c(b)=K+bW$ with $b\in\Theta=[\thl,\thu]$, $\thl\ge 0$.

\begin{definition}[Shortcut of a label]\label{def:shortcut}
Let $P(L)$ be the set of orders picked up in $L$ and $\emptyset\neq T\subseteq P(L)$.
$L\setminus T$ is the subsequence of $L$ without the stops of $T$; $v'$ is its last point
($v'=v_0$ if it is empty), $r'$ its ready time at $v'$ ($r'=t_0$ if empty), and $D'$ its
length. Write $D=D(L)$, $v_m$ and $r=r_m$ for the label itself. Let $\mathcal F(L)$ be the set
of pickups $p_w$ of orders $w\notin P(L)$ with $r+\tau(v_m,p_w)\le l_{p_w}$, taken empty when
$|P(L)|=B$. Define
\begin{align*}
\Delta D(L,T)&:=D-D'-d(v',v_m), &
\delta(L,T)&:=r-r'-\tau(v',v_m),\\
A(L,T)&:=\max\Bigl\{0,\ \max_{p_w\in\mathcal F(L)}\bigl(e_{p_w}-r'-\tau(v',p_w)\bigr)\Bigr\}, &
\beta(L,T)&:=\kappa\,\Delta D+\thl\,\frac{\max\{0,\delta-A\}}{60}.
\end{align*}
\end{definition}

$\Delta D$ and $\delta$ are the distance and time that the label spends on $T$, measured on the
part of the route that is already known; the terms $d(v',v_m)$ and $\tau(v',v_m)$ correct for the
fact that the continuation may start from the last stop of $L$ (this matters when that stop
belongs to $T$). $A$ bounds how much of the time saving later waiting can absorb.

\paragraph{The rule.} \emph{Discard the label $L$ and every extension of it if
$\beta(L,T)\ge q(T)$ for some $\emptyset\ne T\subseteq P(L)$.}

%=====================================================================
\section{Two lemmas on paths}
%=====================================================================

\begin{lemma}[Removing stops never hurts]\label{lem:remove}
Let $\sigma'$ be a subsequence of a stop sequence $\sigma$, both started at $(v_0,t_0)$. Then the
length of $\sigma'$ is at most that of $\sigma$, every retained stop has a service start in
$\sigma'$ no later than in $\sigma$, and the completion time (and, for an occasional driver, the
arrival time at the destination) does not increase.
\end{lemma}

\begin{proof}
This is the argument of [LPF, Lemma~3], which uses only the triangle inequality, nonnegative
service and waiting, and the recursion for $a_k$; feasibility of $\sigma'$ is not needed.
\end{proof}

The next lemma is the new ingredient. It quantifies how a time advantage survives a common
continuation when waiting is allowed.

\begin{lemma}[Absorption of a time advantage]\label{lem:absorb}
Two paths continue with the same stops $w_1,\dots,w_k$. Path~1 leaves point $x$ at ready time
$\eta$, path~2 leaves $x'$ at ready time $\eta'$. Let
$T_1=\eta'+\tau(x',w_1)$ and $T_{j+1}=T_j+s_{w_j}+\tau(w_j,w_{j+1})$ be the no-wait arrival
times of path~2, and suppose $\eta+\tau(x,w_1)\ge T_1+\delta$. Then for every $j$
\[
a^{(1)}_j-a^{(2)}_j\;\ge\;\delta-\max_{l\le j}\,\bigl(e_{w_l}-T_l\bigr)^+ ,
\]
where $a^{(1)}_j,a^{(2)}_j$ are the service starts at $w_j$, and $e_{w}=-\infty$ for stops without
opening time.
\end{lemma}

\begin{proof}
Write $\mathrm{dur}(l\to j):=\sum_{h=l}^{j-1}(s_{w_h}+\tau(w_h,w_{h+1}))$. Unrolling the recursion
gives, for any path entering $w_1$ at no-wait arrival $U_1$,
$a_j=\max\{U_1+\mathrm{dur}(1\to j),\,M_j\}$ with $M_j:=\max_{l\le j}\{e_{w_l}+\mathrm{dur}(l\to j)\}$,
and $M_j$ is the same for both paths. For path~2, $U_1=T_1$ and
$T_1+\mathrm{dur}(1\to j)=T_j$; for path~1, $U_1\ge T_1+\delta$. Hence
$a^{(1)}_j\ge\max\{T_j+\delta,M_j\}$ and $a^{(2)}_j=\max\{T_j,M_j\}$. If $M_j\le T_j$ the difference
is at least $\delta$; if $T_j<M_j\le T_j+\delta$ it is at least $T_j+\delta-M_j$; otherwise it is
at least $0$. In all cases it is at least $\delta-(M_j-T_j)^+$, and
$M_j-T_j=\max_{l\le j}(e_{w_l}-T_l)$ because $T_j=T_l+\mathrm{dur}(l\to j)$.
\end{proof}

%=====================================================================
\section{Main results}
%=====================================================================

\begin{theorem}[Label bound]\label{thm:bound}
Let $L$ be a label, $\emptyset\ne T\subseteq P(L)$, and $\rho$ any \emph{feasible} route of driver
$i$ with prefix $L$ (possibly $\rho=L$). Let $\rho\setminus T$ be $\rho$ without the stops of $T$. Then
\[
K_\rho-K_{\rho\setminus T}\ \ge\ \kappa\,\Delta D(L,T),\qquad
W_\rho-W_{\rho\setminus T}\ \ge\ \frac{\max\{0,\ \delta(L,T)-A(L,T)\}}{60},
\]
and therefore $c_\rho(b)-c_{\rho\setminus T}(b)\ge\beta(L,T)$ for every $b\ge\thl$.
\end{theorem}

\begin{proof}
Write $\rho=L\cdot\sigma$ and consider the intermediate path $\pi:=(L\setminus T)\cdot\sigma$.
Deleting from $\pi$ the stops of $T$ that lie in $\sigma$ (deliveries of orders of $T$ still on
board at the end of $L$) gives $\rho\setminus T$, so by Lemma~\ref{lem:remove}
$K_\pi\ge K_{\rho\setminus T}$ and $W_\pi\ge W_{\rho\setminus T}$. It is therefore enough to
compare $\rho$ with $\pi$; also $K_\rho\ge K_\pi$ and $W_\rho\ge W_\pi$ by the same lemma.

\emph{Distance.} If $\sigma$ is empty (a gigworker route ending at $v_m$), $D(\rho)-D(\pi)=D-D'\ge\Delta D$.
Otherwise let $u$ be the first point of $\sigma$ (for an occasional driver possibly the
destination). Then $D(\rho)-D(\pi)=D+d(v_m,u)-D'-d(v',u)\ge D-D'-d(v',v_m)$ by
$d(v',u)\le d(v',v_m)+d(v_m,u)$. For an occasional driver the direct length cancels, and both
detours are nonnegative, so the positive part in $K$ does not interfere.

\emph{Time.} If $\sigma$ is empty, $W_\rho-W_\pi=(r-r')/60\ge\delta/60$ since $\tau\ge0$. Otherwise
apply Lemma~\ref{lem:absorb} with $x=v_m$, $\eta=r$, $x'=v'$, $\eta'=r'$: the hypothesis holds
because $r+\tau(v_m,w_1)\ge r'+\tau(v',w_1)+\delta$ follows from
$\tau(v',w_1)\le\tau(v',v_m)+\tau(v_m,w_1)$. At the last stop (and at the destination), the
completion times therefore differ by at least $\delta-\max_l(e_{w_l}-T_l)^+$. A stop $w_l$ of
$\sigma$ with an opening time is a pickup of an order not picked in $L$; it can occur only if
$|P(L)|<B$, and feasibility of $\rho$ requires $r+\tau(v_m,w_l)\le l_{w_l}$ (triangle
inequality), so $w_l\in\mathcal F(L)$. Since $T_l\ge r'+\tau(v',w_l)$ (triangle inequality,
nonnegative service), the maximum is at most $A(L,T)$. Combined with $W_\rho\ge W_\pi$ this
gives the time bound.

\emph{Cost.} $c_\rho(b)-c_{\rho\setminus T}(b)=(K_\rho-K_{\rho\setminus T})+b\,(W_\rho-W_{\rho\setminus T})$
with a nonnegative slope, so it is at least its value at $\thl$, which is at least $\beta$.
\end{proof}

\begin{theorem}[The rule preserves the local frontier exactly]\label{thm:preserve}
Let $R_i$ be the pool of Algorithm~A (complete enumeration, per-bundle Pareto filtering, one
representative per pair $(K,W)$), and $R'_i$ the pool obtained when Algorithm~A in addition
discards some or all labels satisfying the rule, together with their extensions. Under the
assumptions of [LPF, Lemma~3] and $\thl<\thu$,
\[
\Ks_i(R'_i)=\Ks_i(R_i).
\]
Consequently the pipeline ``Algorithm~A with the rule, then $\Pi^\star$'' outputs exactly the
same pool as before, and [LPF, Theorem~4 and Corollary~6] apply verbatim: the WDP value, every
removal value and the Clarke-pivot payments are unchanged.
\end{theorem}

\begin{proof}
Let $F$ be the set of all feasible routes, so $R_i$ is the Pareto representative set of $F$
and $R'_i$ that of $F$ minus the discarded routes.

\emph{(i) Discarded routes of $R_i$ lie outside $\Ks_i(R_i)$.} Let $\rho\in R_i$ extend a discarded
label with witness $T$. By Theorem~\ref{thm:bound},
$c_\rho(b)\ge c_{\rho\setminus T}(b)+\beta\ge c_{\rho\setminus T}(b)+q(T)$ for all $b\in\Theta$.
The route $\rho\setminus T$ is feasible ([LPF, Lemma~3]) and by downward closure there is, for
every $b$, some $\hat\rho\in R_i\cup\{\0_i\}$ with bundle $S_\rho\setminus T$ and
$c_{\hat\rho}(b)\le c_{\rho\setminus T}(b)$. Since $\hat\rho\in D(\rho)$,
$E_\rho(b)\le c_{\hat\rho}(b)+q(T)\le c_\rho(b)$; hence $\mu_\rho\le0$ on $\Theta$ and
$\rho\notin\Ks_i(R_i)$.

\emph{(ii) $\Ks_i(R_i)\subseteq R'_i$.} A route of $\Ks_i(R_i)$ is not discarded by (i), and it is
Pareto-undominated in $F$, hence in any subset of $F$ containing it.

\emph{(iii) $\Ks_i(R_i)\subseteq\Ks_i(R'_i)$.} Take $r\in\Ks_i(R_i)$. Compared with $D_{R_i}(r)$, the
set $D_{R'_i}(r)$ lacks some routes (which can only raise $E_r$) and may contain routes $r_2$
that were Pareto-dominated in $F$; such an $r_2$ is dominated by a route $r_1\in R_i$ with the same
bundle, so $c_{r_2}+q(\cdot)\ge c_{r_1}+q(\cdot)$, which is at least $E^{R_i}_r$ when $r_1\ne r$; when $r_1=r$,
Pareto dominance gives $c_{r_2}(b)>c_r(b)$ for every $b>0$. Hence on the nonempty open set where
$\mu^{R_i}_r>0$ (minus possibly $b=0$), $E^{R'_i}_r>c_r$, and $r\in\Ks_i(R'_i)$.

\emph{(iv) $\Ks_i(R'_i)\subseteq\Ks_i(R_i)$.} Let $r\in R'_i$ with $r\notin\Ks_i(R_i)$. If $r\in R_i$,
[LPF, Lemma~5] gives, for each $b$, some $u\in\Ks_i(R_i)\cup\{\0_i\}$ with $S_u\subseteq S_r$ and
$c_u(b)+q(S_r\setminus S_u)\le c_r(b)$; $u\in R'_i$ by (ii), so $\mu^{R'_i}_r\le 0$. If
$r\notin R_i$, $r$ is Pareto-dominated in $F$ by some $r_1\in R_i$ with the same bundle; $r_1$ must
have been discarded (otherwise $r$ would not survive the Pareto filter of $R'_i$), so
$r_1\notin\Ks_i(R_i)$ by (i), and the same lemma applied to $r_1$ yields $u\in R'_i$ with
$c_u+q(S_r\setminus S_u)\le c_{r_1}\le c_r$. In both cases $r\notin\Ks_i(R'_i)$.
\end{proof}

\begin{remark}[What the rule is and is not]
The rule reads no bid (only the public constant $\thl$ and FD prices) and uses only driver $i$'s
own data, so it is a local, bid-independent rule in the sense of [LPF]. By [LPF, Theorem~4(c)],
\emph{no} rule of this kind can discard a route of $\Ks_i$; the best any such label rule can do
is to cut every label whose subtree contains no route of $\Ks_i$. Section~\ref{sec:evidence}
measures how close the rule comes to that ceiling.
\end{remark}

\begin{remark}[Combination with label dominance]
If Algorithm~A also deletes a label $L_2$ because another label $L_1$ with the same end stop and
the same picked/delivered sets dominates it (ready time, length and cost no larger), every
completion of $L_2$ is dominated by the corresponding completion of $L_1$: the recursion for $a_k$
is nondecreasing in the starting time, so the same continuation from $L_1$ is feasible, no longer
and finishes no later. If $L_1$ is discarded by the rule, its completions lie outside
$\Ks_i$, and so do those of $L_2$, since a same-bundle route that costs no less for every $b$
cannot be in $\Ks_i$. The proof of Theorem~\ref{thm:preserve} is unchanged.
\end{remark}

\begin{remark}[Delivery time windows with an opening time]\label{rem:dlb}
If deliveries also have opening times, include in $\mathcal F(L)$ the deliveries of orders on
board at the end of $L$ and of every order that may still be picked up. Lemma~\ref{lem:absorb}
is unchanged.
\end{remark}

%=====================================================================
\section{Why each term is needed: a counterexample and a negative result}
%=====================================================================

\begin{example}[Direct transfer of rollback pruning is unsafe]\label{ex:counter}
Rollback pruning compares a path with the same path skipping a node and prunes if skipping is
not worse; in the ESPPRC this is safe because costs are additive over arcs and time is only a
resource, so arriving earlier never costs anything. Here working time is part of the cost and
waiting is allowed. Transferring rollback directly amounts to setting $A\equiv 0$ in
Definition~\ref{def:shortcut}. On a synthetic instance (seed~2, gigworker $g_0$, $B=3$), for the
label $L=(p_{o_2},d_{o_2})$ and $T=\{o_2\}$ with $q_{o_2}=17.13$, the uncorrected bound credits a
time saving of $42.9$ min, giving $17.72\ge q_{o_2}$, and the label is discarded. In the feasible
completion $(p_{o_2},d_{o_2},p_{o_0},d_{o_0})$, however, the pickup of $o_0$ opens late: without
$o_2$ the driver would arrive early and wait, so the true time saving is only $7.8$ min and the
true saving at $b=\thl$ is $7.24<17.13$. Serving $o_2$ in this route is therefore worth more than
delegating it to FD, and the uncorrected rule deletes a route it must keep. With the correction
$A$ the label is kept.
\end{example}

The mutation tests in Section~\ref{sec:evidence} show the same at scale for both correction
terms.

\paragraph{Negative result: a forward lower-bound matrix does not work.}
The most natural first idea is to decide \emph{before} visiting $p_o$: from the current stop $v$,
bound the extra distance by
$\min_x\{d(v,p_o)+d(p_o,x)-d(v,x)\}$ over all possible next stops $x$ (a matrix that can be
precomputed), and the extra time analogously. This rule is safe (the same lemmas apply) but on
the synthetic instances it never fired. The median distance credit was $0.5$--$0.75$ USD against
FD prices of $15$--$21$ USD, because some other stop almost always lies close to the direction of
$p_o$; and the median time credit was $0$, because the detour (about $5$ min) is fully absorbed
by the waiting for later-released pickups (median absorption $34$--$46$ min). Both reasons are
structural: before $p_o$ is visited, the continuation is unknown, and the bound must hold for
the cheapest one. The rule of this note is the same idea evaluated one stop \emph{after} $p_o$,
when the detour to $p_o$ is known exactly; the correction terms are what make that evaluation
safe for every continuation.

%=====================================================================
\section{Implementation}
%=====================================================================

Each label carries, for every picked order $j$, the triple $(D'_j,v'_j,r'_j)$ of its shortcut
$L\setminus\{j\}$, as an additional resource. Extending the label by a stop $u$ updates each triple
in $O(1)$: if $u$ is a stop of $j$ the triple is unchanged; otherwise the shortcut travels to $u$
exactly like the label. When $j$ is picked up, its triple is the parent label's state. The check
first evaluates $\beta$ with $A=0$, which is an upper bound; only if that reaches $q_j$ is $A$
computed. The per-extension overhead is therefore $O(|P(L)|)\le O(B)$ plus rare $O(n)$
evaluations of $A$. Subsets $T$ with $|T|\ge2$ added almost nothing in the experiments and are
not needed.

%=====================================================================
\section{Evidence}\label{sec:evidence}
%=====================================================================

All numbers come from the independent brute-force reference implementation of [LPF] (synthetic
Euclidean instances, $n=8$ orders, two gigworkers and two occasional drivers, $\Theta=[18,25]$).
They check the mathematics; they are not thesis results.

\begin{center}
\begin{tabular}{@{}lll@{}}
\toprule
Check & $B=3$ (3 seeds) & $B=4$ (2 seeds) \\
\midrule
Rule firings audited & 5{,}868 & 33{,}292 \\
Feasible completions checked against the bound & 9{,}975 & 87{,}744 \\
Bound violations & 0 & 0 \\
$\Ks$ identical with and without the rule & yes & yes \\
Max WDP value difference (random bids) & $2.8\times10^{-14}$ & $1.3\times10^{-13}$ \\
\midrule
Mutation $A\equiv0$ (direct rollback transfer) & 1{,}339 violations, $\Ks$ changed & --- \\
Mutation without junction terms $d(v',v_m),\tau(v',v_m)$ & 4{,}173 violations, $\Ks$ changed & --- \\
\bottomrule
\end{tabular}
\end{center}

In both mutations the WDP values on 40 random bid profiles per seed were still identical: the
wrongly deleted routes are rarely optimal, so bid sampling alone would not detect the error. The
completion-level bound audit and the $\Ks$ comparison did.

\paragraph{Distance to the ceiling.} Counting labels that lie on a path to some feasible route:
an ideal local rule (Theorem~4(c) of [LPF]) keeps $2.7\%$ of them at $B=3$ and $0.6\%$ at $B=4$;
the rule keeps $56\%$ and $50\%$, i.e.\ it closes $45\%$ and $50\%$ of the achievable gap.

\paragraph{Work and time.} With the incremental implementation, label extensions fell by
$19$--$44\%$ at $B=3$ and $26$--$61\%$ at $B=4$. Wall-clock time in this plain enumerator ranged
from $0.8\times$ to $2\times$ the baseline. The decisive factor is the FD price level relative
to route costs (median over three seeds and two gigworkers, $B=4$):

\begin{center}
\begin{tabular}{@{}lccc@{}}
\toprule
FD prices scaled by & $0.6$ & $1.0$ & $1.4$\\
\midrule
Label extensions saved & $73$--$76\%$ & $27$--$30\%$ & $7$--$9\%$\\
Wall-clock speed-up & $2.8$--$3.1\times$ & $\approx1.0\times$ & $0.75\times$\\
\bottomrule
\end{tabular}
\end{center}

Compressing release times by a factor of four changed these figures by at most three points. On
the RQ1 instances $\Pi^\star$ removes $77$--$97\%$ of the pool, against about $61\%$ on these synthetic
instances, which suggests that FD is relatively cheaper there and that the rule is in its
effective regime; this must be measured, not assumed.

%=====================================================================
\section{Position in the literature}
%=====================================================================

The principle --- compare a partial path with itself minus a node and prune if the node is not
worth its price --- is the rollback pruning of Lozano, Duque \& Medaglia (2016,
\emph{Transportation Science}, \S4.3), built on the dominance relation of Feillet et al.\ (2004);
there the price is the node's dual value. It must be cited as the origin of the rule. The
adaptation adds four things, each needed for correctness or for the mechanism: (a) the
correction $A$ for waiting absorption, without which the rule is unsafe when time enters the
objective (Example~\ref{ex:counter}); (b) removal of pickup--delivery \emph{pairs}, including
orders still on board and a removed last stop (junction terms); (c) comparison uniformly over
$\Theta$ with a bid-independent price, so that the pool stays fixed before bids; (d) the exact
guarantee of Theorem~\ref{thm:preserve} that the local frontier, and hence the VCG outcome, is
unchanged. Label-dominance rules for the PDPTW (Dumas, Desrosiers \& Soumis, 1991; Ropke \&
Cordeau, 2009) and bounding strategies of the pulse family have not yet been checked for a closer
precedent; this must be done before (a)--(d) are claimed as new.

\section*{References to verify before use}
Lozano, L., Duque, D.\ \& Medaglia, A.L.\ (2016), An exact algorithm for the elementary
shortest path problem with resource constraints, \emph{Transportation Science} (rollback
pruning, \S4.3; read directly from the project copy, online version dated 2015).
Feillet, D., Dejax, P., Gendreau, M.\ \& Gueguen, C.\ (2004), \emph{Networks} (dominance for the
ESPPRC; not read directly). Dumas, Desrosiers \& Soumis (1991), \emph{EJOR}. Ropke \& Cordeau
(2009), \emph{Transportation Science}.

\end{document}
